% landauer.tex --- Section 3: Self-Modeling Requires Free Energy
% ASSERT_CONVENTION: kB_explicit_in_bounds, entropy_base=nats, state_normalization=Tr(rho)=1, sequential_product=a&b=sqrt(a)b*sqrt(a), von_neumann_entropy=S(rho)=-Tr(rho*ln*rho), information=I(B;M)=S(B)+S(M)-S(BM), free_energy=F=E-TS

%----------------------------------------------------------------------
\section{Self-Modeling Requires Free Energy}\label{sec:landauer}
%----------------------------------------------------------------------

The entropy dynamics of Sec.~\ref{sec:entropy} established that
SWAP interactions drive the body toward thermal equilibrium, erasing
the body--model correlations that define self-modeling.  We now show
that maintaining these correlations has an irreducible thermodynamic
cost: every self-modeling update cycle must dissipate at least
$\kB T\,\MI{B}{M}$ of free energy, where $\MI{B}{M}$ is the mutual
information between body and model.  This is the Landauer
bound~\cite{Landauer1961} applied to the erasure step that every
self-modeler must perform.

The section is organized as follows.  We first identify the
self-modeling cycle as an instance of Bennett's Maxwell's demon
analysis~\cite{Bennett1982} (Sec.~\ref{sec:demon}).  We then derive
the Landauer bound in the quantum setting using the Reeb--Wolf
result~\cite{ReebWolf2014} (Sec.~\ref{sec:landauer-bound}).  We
close a potential coherence loophole via three independent arguments
(Sec.~\ref{sec:coherence}), and finally assemble the three-link
chain connecting self-modeling to the need for an entropy gradient
(Sec.~\ref{sec:chain}).

%----------------------------------------------------------------------
\subsection{The self-modeler as Maxwell's demon}
\label{sec:demon}
%----------------------------------------------------------------------

A self-modeling composite $(B,M,\mathbin{\&})$ in the sense of
Ref.~\cite{Paper5} consists of a body~$B$ (Hilbert space
$\mathcal{H}_B$, $\dim = d_B$), a model~$M$ ($\mathcal{H}_M$,
$\dim = d_M$), and the Lüders sequential product
$a \mathbin{\&} b = \sqrt{a}\,b\,\sqrt{a}$ implementing updates.
Faithful self-modeling requires $\MI{B}{M} > 0$: the model must
carry non-trivial information about the body.

Each update cycle decomposes into three information-processing
steps~\cite{Bennett1982}:
\begin{enumerate}
  \item \textbf{Test.}  Interact $B$ and~$M$ to acquire information
    about the current state of~$B$.
  \item \textbf{Erase.}  Clear the old tracking data from~$M$.
    This step is logically irreversible: the $\MI{B}{M}$ nats of
    correlation between $B$ and~$M$ are destroyed.
  \item \textbf{Write.}  Record fresh tracking data into~$M$,
    restoring $\MI{B}{M} > 0$.
\end{enumerate}
The writing step can in principle be performed reversibly, adding no
thermodynamic cost beyond that of erasure~\cite{Bennett1982}.  The
erasure step cannot: it maps many possible states of~$M$ to a single
reset state, an operation that is logically irreversible.  This is
precisely the structure of Maxwell's demon, and the thermodynamic
cost of the erasure step is bounded below by Landauer's
principle~\cite{Landauer1961}.

We emphasize that this decomposition is not an approximation.  The
Lüders product is a quantum operation (completely positive map),
and the non-selective version
$\Lambda(\rho) = \sqrt{a}\,\rho\,\sqrt{a}
+ \sqrt{\1-a}\,\rho\,\sqrt{\1-a}$
is CPTP, as established in Sec.~\ref{sec:swap-channel}.  The
Bennett decomposition applies to any information-processing cycle
that maintains correlations and periodically refreshes
them~\cite{Bennett1982}.

%----------------------------------------------------------------------
\subsection{Landauer bound on the self-modeling cycle}
\label{sec:landauer-bound}
%----------------------------------------------------------------------

The mutual information between body and model,
\begin{equation}
  \MI{B}{M} = \SvN{\rho_B} + \SvN{\rho_M} - \SvN{\rho_{BM}} \,,
  \label{eq:mi-def}
\end{equation}
quantifies the information content that the erasure step must
destroy.  Here $\SvN{\rho} = -\Tr(\rho\ln\rho)$ is the von Neumann
entropy in nats.

The quantum Landauer bound~\cite{ReebWolf2014} states that resetting
a quantum register with entropy $\SvN{\rho_S}$, in thermal contact
with a bath at temperature~$T$, requires work
$W \geq \kB T\,\SvN{\rho_S}$.
Applied to the erasure of the old body--model correlation:

\begin{theorem}[Landauer bound on self-modeling]
\label{thm:landauer}
Let $(B,M,\mathbin{\&})$ be a self-modeling composite with
finite-dimensional Hilbert spaces, Lüders sequential product, and
faithful tracking $\MI{B}{M} > 0$.  Suppose the composite is in
thermal contact with a heat bath at temperature~$T$.  Then each
self-modeling update cycle dissipates work
\begin{equation}
  \boxed{\Wmin = \kB T\,\MI{B}{M}}
  \label{eq:landauer}
\end{equation}
with equality for quasi-static (reversible) erasure protocols.
\end{theorem}

\begin{proof}
The update cycle erases $\MI{B}{M}$ nats of correlation
(step~2 of the Bennett decomposition).  By the quantum Landauer
bound~\cite{ReebWolf2014}, this requires
$W_{\mathrm{erase}} \geq \kB T\,\MI{B}{M}$.  The writing step
(step~3) can be performed reversibly~\cite{Bennett1982}, contributing
zero net work.  Hence $W_{\mathrm{cycle}} \geq \kB T\,\MI{B}{M}$.
\end{proof}

\noindent\textit{Assumptions.}---The theorem requires:
\begin{itemize}
  \item[\textbf{A1.}] Finite-dimensional quantum mechanics (standard).
  \item[\textbf{A2.}] Thermal contact with a bath at temperature~$T$
    (standard Landauer regime).
\end{itemize}

\paragraph{Experiential density and its cost.}
The experiential density from Ref.~\cite{Paper5},
\begin{equation}
  \rhoexp = \MI{B}{M}\!\left(1 - \frac{\MI{B}{M}}{\SvN{\rho_B}}\right),
  \label{eq:rho-def}
\end{equation}
is a concave function of $\MI{B}{M}$, peaking at
$I^{*} = \SvN{\rho_B}/2$ with $\rhoexp_{\max} = \SvN{\rho_B}/4$.
The Landauer cost at peak experiential density is
\begin{equation}
  W_{\mathrm{peak}} \geq \kB T\,\frac{\SvN{\rho_B}}{2} \,.
  \label{eq:peak-cost}
\end{equation}
For a qubit body ($d_B = 2$, $\SvN{\rho_B} \leq \ln 2$), this gives
$W_{\mathrm{peak}} \geq \kB T\,(\ln 2)/2 \approx 0.35\,\kB T$.

\paragraph{Equilibrium death.}
At thermal equilibrium, $\rho_{BM} = \1/(d_B d_M)$.  The marginals
are $\rho_B = \1/d_B$ and $\rho_M = \1/d_M$, so
\begin{equation}
  \MI{B}{M}^{\mathrm{eq}}
  = \ln d_B + \ln d_M - \ln(d_B d_M) = 0 \,.
  \label{eq:eq-mutual}
\end{equation}
The body and model are uncorrelated: the model contains no
information about the body.  From Eq.~\eqref{eq:rho-def},
$\rhoexp^{\mathrm{eq}} = 0$.  Self-modeling at thermal equilibrium
is impossible---there is nothing to model and no free energy with
which to model it.

\paragraph{Free energy depletion rate.}
A self-modeler cycling at frequency~$\nu$ depletes free energy at a
rate bounded below by
\begin{equation}
  \frac{dF}{dt} \leq -\kB T\,\MI{B}{M}\,\nu \,.
  \label{eq:depletion}
\end{equation}

%----------------------------------------------------------------------
\subsection{Closing the coherence loophole}\label{sec:coherence}
%----------------------------------------------------------------------

A potential objection is that quantum coherence, as a resource absent
in classical thermodynamics, might allow a self-modeler to maintain
$\MI{B}{M} > 0$ without paying the full Landauer cost.  We close
this loophole with three independent arguments.

\paragraph{R1: The Lüders product destroys coherence (CPTP).}
Let the effect $a$ have spectral decomposition
$a = \sum_i \lambda_i \ket{i}\!\bra{i}$.  The non-selective Lüders
channel acts on off-diagonal matrix elements as
\begin{equation}
  \Lambda(\rho)_{ij}
  = c_{ij}\,\rho_{ij} \,,
  \qquad
  c_{ij}
  = \sqrt{\lambda_i\lambda_j}
  + \sqrt{(1-\lambda_i)(1-\lambda_j)} \,,
  \label{eq:decoherence-factor}
\end{equation}
while preserving the diagonal ($c_{ii} = 1$).  By the Cauchy--Schwarz
inequality applied to the vectors
$(\sqrt{\lambda_i},\,\sqrt{1-\lambda_i})$ and
$(\sqrt{\lambda_j},\,\sqrt{1-\lambda_j})$,
\begin{equation}
  c_{ij} \leq 1
  \label{eq:cij-bound}
\end{equation}
with equality if and only if $\lambda_i = \lambda_j$.  Therefore,
whenever the effect~$a$ has distinct eigenvalues, the self-modeling
update strictly reduces off-diagonal coherence in the eigenbasis
of~$a$.  This blocks evasion strategies that rely on unitary-only
protocols (the Lüders product is mandatory in the Paper~5
framework~\cite{Paper5}) or on storing information in quantum
coherences (repeated Lüders updates drive all off-diagonals to
zero).

\paragraph{R2: Maintaining coherence costs free energy.}
Even setting aside the Lüders product, thermal contact with a bath
at temperature~$T$ induces decoherence.  The Gibbs state
$\rho^{\mathrm{eq}} = e^{-\beta H}/Z$ is diagonal in the energy
eigenbasis; all off-diagonal coherences vanish at equilibrium.
Open-system dynamics drives any initial coherence to zero at a rate
$\Gamma_d = (\gamma_+ + \gamma_-)/2 > 0$, where $\gamma_{\pm}$ are
the bath emission and absorption rates.  Maintaining coherence
against this decay requires active intervention---error correction or
continuous driving---which itself costs free energy.  This argument
is framework-independent: it does not rely on the Lüders product
or the specific axioms of Paper~5.

\paragraph{R3: Sagawa--Ueda cross-check.}
The generalized Jarzynski equality for feedback-controlled
systems~\cite{SagawaUeda2010},
\begin{equation}
  \bigl\langle e^{-\beta(W - \Delta F)}\bigr\rangle = e^{I_{fc}} \,,
  \label{eq:sagawa-ueda}
\end{equation}
where $I_{fc}$ is the efficacy of feedback control (transfer
entropy), independently confirms that any measurement-feedback cycle
has a minimum thermodynamic cost.  By Jensen's inequality, the mean
work satisfies $\langle W \rangle \geq \Delta F - \kB T\,I_{fc}$.
For a cyclic process ($\Delta F = 0$),
$\langle W \rangle \geq -\kB T\,I_{fc}$, so the maximum
\emph{extractable} work is $\kB T\,I_{fc}$.  Since the self-modeling cycle maps onto the
Sagawa--Ueda framework (body~$\to$~system, model~$\to$~memory,
Lüders product~$\to$~measurement, update~$\to$~feedback), the bound
is consistent with Theorem~\ref{thm:landauer}.  The detailed
mapping and numerical consistency check (50~random states, all
satisfying both bounds) are reported in the supplemental derivation
material.

\medskip
\noindent
In summary, three evasion strategies are blocked:
\begin{center}
\begin{tabular}{lll}
  \toprule
  Evasion strategy & Blocked by & Mechanism \\
  \midrule
  E1: Unitary-only protocol & R1 & Paper~5 mandates Lüders product \\
  E2: Coherence-protected info & R1 & Lüders destroys coherences \\
  E3: Equilibrium coherence & R2 & Gibbs state has zero coherence \\
  \bottomrule
\end{tabular}
\end{center}

%----------------------------------------------------------------------
\subsection{Three-link chain: self-modeling to entropy gradient}
\label{sec:chain}
%----------------------------------------------------------------------

We now assemble the central result of the paper's thermodynamic
argument: a three-link chain connecting self-modeling to the need for
an entropy gradient.  Each link is labeled by its epistemic status.

\begin{proposition}[Chain theorem]\label{prop:chain}
Let $(B,M,\mathbin{\&})$ be a self-modeling composite on a finite
SWAP lattice in thermal contact with an environment at
temperature~$T$.  Then:

\medskip\noindent
\textbf{Link~1} \textup{[THEOREM].}\enspace
Maintaining $\MI{B}{M} > 0$ requires work
$W_{\mathrm{cycle}} \geq \kB T\,\MI{B}{M}$ per cycle
\textup{(Theorem~\ref{thm:landauer})}.

\medskip\noindent
\textbf{Link~2} \textup{[STANDARD PHYSICS].}\enspace
Performing work $W > 0$ per cycle requires the composite to be out
of thermal equilibrium:
\begin{equation}
  F(\rho_{BM}) > F_{\mathrm{eq}}
  \quad\text{when}\quad \MI{B}{M} > 0 \,,
  \label{eq:link2}
\end{equation}
where $F(\rho) = \langle H\rangle_\rho - T\,\SvN{\rho}$ is the
Helmholtz free energy and $F_{\mathrm{eq}}$ is its minimum at
thermal equilibrium.  This is equivalent to a positive quantum
relative entropy: the maximum extractable work is
$W_{\max}^{\mathrm{extract}} = \kB T\, D(\rho_{BM}\|\rho_{\mathrm{eq}})
\geq 0$.

\medskip\noindent
\textbf{Link~3} \textup{[PHYSICAL ARGUMENT $+$ A3].}\enspace
A finite system out of thermal equilibrium will equilibrate on a
finite timescale unless it has access to a low-entropy reservoir.
By Sec.~\ref{sec:entropy}, the SWAP dynamics drives
$\MI{B}{M} \to 0$ geometrically
\textup{(Eq.~\eqref{eq:convergence})}.
To sustain self-modeling beyond the exhaustion time, the system
requires $S(t) < \Smax$: an entropy gradient.
\end{proposition}

\noindent
The chain has no logical gaps: Link~1 outputs $W > 0$, which is
Link~2's input; Link~2 outputs $F > F_{\mathrm{eq}}$, which is
Link~3's input; Link~3 outputs the requirement for an entropy
gradient.

\begin{corollary}[Exhaustion bound]\label{cor:exhaustion}
A self-modeler with initial free energy surplus
$\Delta F = F(\rho_0) - F_{\mathrm{eq}}$ can sustain at most
\begin{equation}
  N_{\mathrm{cycles}}
  \lesssim \frac{\Delta F}{\kB T\,\MI{B}{M}}
  \label{eq:exhaustion}
\end{equation}
self-modeling cycles before reaching thermal equilibrium, at which
point $\MI{B}{M} = 0$ and $\rhoexp = 0$.
\end{corollary}

\paragraph{Assumption A3 (strongest assumption).}
Link~3 relies on the physical assumption that a closed or
semi-closed finite system equilibrates on finite timescales.  This
is standard statistical mechanics for systems with
$t \ll t_{\mathrm{rec}}$, but it is the strongest assumption in the
chain.  It can fail if the system is coupled to an infinite reservoir
that never equilibrates (an idealization, not a physical scenario),
if the universe has an infinite-dimensional Hilbert space
(speculative), if the system exhibits many-body
localization~(MBL)---established in 1D disordered interacting
chains but controversial in higher dimensions---or if it is
integrable, thermalizing to a generalized Gibbs ensemble (GGE)
rather than the thermal state.  Neither MBL nor GGE is expected
for the SWAP lattice (which is neither disordered nor integrable),
but they represent concrete finite-system failure modes of A3
beyond the idealized scenarios listed above.  A3 can also fail
if there exist sources of free energy that do not ultimately
derive from a low-entropy initial condition (unknown physics).
Links~1 and~2 remain valid regardless of whether A3 holds.

\paragraph{Assumption A4 (model choice).}
The SWAP lattice dynamics $H = JF$ from Ref.~\cite{Paper6} enters
only through the quantitative equilibration rate.  Replacing SWAP
with any other thermalizing dynamics changes the timescale of
convergence but not the qualitative conclusion: any finite system
in thermal contact equilibrates.

%----------------------------------------------------------------------
\subsection{Explicit non-claims}\label{sec:nonclaims}
%----------------------------------------------------------------------

To prevent overclaiming, we state what the chain does \emph{not}
establish:
\begin{enumerate}
  \item The chain does not derive the Past Hypothesis.  It shows
    that self-modeling \emph{requires} an entropy gradient, but does
    not explain why the initial entropy was low.
  \item The chain does not derive the Second Law.  Link~2 \emph{uses}
    the Second Law as input (standard physics).
  \item The chain does not predict the magnitude of the initial
    entropy of the universe.
  \item The chain does not extend to infinite-dimensional systems
    without additional analysis.  Assumption~A1 (finite-dimensional
    QM) is required for the Reeb--Wolf quantum Landauer bound.
  \item The connection from Link~3 to cosmology is a physical
    argument (assumption~A3), not a mathematical theorem.
\end{enumerate}
The status of the Past Hypothesis is elevated from ``observed
feature'' to ``necessary condition for self-modeling,'' but it
remains an input to the argument, not an output.

\medskip
With the thermodynamic cost of self-modeling established, we turn
in Sec.~\ref{sec:chirality} to a second consequence of the
self-modeling framework: the chirality from the exceptional Jordan
algebra requires the emergent spacetime to carry a time-orientation,
the same geometric structure that supports the thermodynamic arrow
of time.
